\scanpage{96}
jouant un rôle universel pour \deleted{\unsure{ces}} \inserted{de telles} données ($\underline I$ torseur sous $(\mathbb Z/3\mathbb Z)_S$, $t\in\Gamma(\Sigma^*_{\underline I}/S)$, $\omega\in\Gamma(S,\underline\omega_{\underline I})$ correspondant à la structure de torseur sur $\underline I$, et donnant lieu à une trivialisation de $\underline\omega_{\underline I}$ sur $S$ \ldots). Les deux pts de vue sont équivalents, quand on se fixe une section $i$ de $\mu_4^\times$, i.e. quand on travaille au dessus de la base $\widetilde S_0=\mu_{4\,S_0}^*$ ($S_0=\operatorname{Spec}\mathbb Z[1/2]$).

\begin{center}---\end{center}

Explicitons entièrement la notion d'épinglage de Legendre, du point de vue calculatoire, sur un corps de base alg. clos $k$ pour simplifier, en décrivant la courbe elliptique $E$ sur $k$ en termes de
\[
\boxed{I\bigl(={}_2E(k)^*\bigr)}\quad\text{(ens à trois éléments), du} \tag*{(58)}
\]
$k$-vectoriel
\[
\boxed{V_I=\deleted{\unsure{k^I|k}}\operatorname{Ker}(k^I\xrightarrow{\text{somme}}k)} \tag*{(59)}
\]
(\unsure{de là|et là|d'où} $\Sigma_I=\mathbb P^1\deleted{\unsure{roj}}(V_I)$, avec \deleted{\unsure{q|\delta}} des $\dot\xi_i\in\Sigma_I(k)$ $(i\in I)$ d'où $\Sigma_I^*=\Sigma_I-\{\dot\xi_i\}_{i\in I}$), et un $t\in\Sigma_I^\times$ i.e. une droite
% slanted
\margin{$\dot\xi_i=k.(e_j-e_k)$ où $(e_i)$ est la base can. de $V_I$}
\[
F_t\subset V_I\,, \tag*{(60)}
\]
enfin un $k$-module inv. $\underline t\deleted{_E}$ \inserted{$(=\underline t_E)$} et un isom
\[
\unsure{\nu|v}(=\unsure{\nu_E|v_E}):\underline t\deleted{\unsure{_E}}^{\otimes2}\simeq T_{\Sigma_I,t}\,. \tag*{(61)}
\]
Il faut d'abord expliciter calculatoirement $T_{\Sigma_I,t}$
% archive sheet number: 95
